\secciongrande{Nuestro primer script javascript}

\resumen{
  Se puede ejecutar una instrucción simple de javascript, como \texttt{document.write("¡Hola Onliner!")}, para mostrar texto directamente en la página web. Es fundamental utilizar el punto y coma (\texttt{;}) para delimitar la terminación de cada instrucción en javascript.
}

\nota{Para el examen}{
¿Qué función cumple el punto y coma (\texttt{;}) al final de una línea de código javascript? Indica al intérprete del navegador que la instrucción ha finalizado.
}

\begin{codigo}
document.write("¡Hola Onliner!");
\end{codigo}



\textbf{Estructura completa del script}

\begin{codigo}
\cmd{<script>}
document.write("¡Hola Onliner!");
\cmd{</script>}
\end{codigo}



\textbf{Nota sobre la terminación}

\nota{Para el examen}{
La etiqueta de cierre \texttt{</script>} indica al navegador que ha terminado de procesar el código javascript y debe volver a esperar código html.
}

\begin{codigo}
\cmd{</script>}
\end{codigo}



]